\rev{The failures of Table~\ref{tab:taxonomy} share one root cause:
the plan relies on the model to compute physical quantities that could
instead be obtained exactly, either from a sensor measurement or from
an action the plan has already committed to. This observation
motivates the paper's main contribution. We extend the pipeline from
two layers to three by inserting \harness{} (\harnesslong{}) between
the high-level planner and the low-level executor. \harness{} is a
deterministic layer: it receives each validated plan before execution
and replaces the numeric parameters most exposed to the identified
failure modes with values it computes itself, using only information
available at plan time.}

The design principle fits in one sentence: \emph{the model decides
what to do and where; \harness{} computes how much.} The model still
chooses which object to pick and which zone to place it in. Geometric
arithmetic is pulled out of the model and into code wherever the
answer is computable from a sensor measurement or from a paired action
already in the plan. The harness is implemented as a small set of
deterministic correction functions; each is described below, together
with the failure it \rev{resolves}.

\subsection{Heights anchored to committed actions and measured depth}
\label{sec:fix-height}

Two corrections govern vertical geometry. The first \rev{resolves} the
phantom top-surface grasp. For every release, \harness{} recomputes
the lift height from the \emph{same} descent the executor actually
commanded during the paired pick, $\height - \descent$ in the notation
of Case study 1. A value derived from the plan's own committed action
cannot disagree with what the arm did; a separately model-computed
value can, and did.

The second \rev{resolves} the double-counted stacking height. For a
stacking target, \rev{\harness{} measures the target's real height
directly: it takes the depth-measured top surface and the table plane,
and uses the vertical distance between the two as the height. The
height is never taken as negative, and no minimum is imposed, so a
flat zone marked on the table correctly reads as having near-zero
height.} It then rewrites the target's base reference so that the
planner's usual reconstruction of the top surface becomes correct. The
measured value is inserted into the representation the planner already
expects, instead of being mixed with the model's own guess.

\subsection{Footprint-based target matching}
\label{sec:fix-footprint}

\rev{To decide whether a planned placement position falls within a
registered target zone, \harness{} checks whether the position lies
inside the zone's real footprint, the rectangle spanned by the zone's
size around its center, rather than measuring the distance to the
zone's center point. This resolves the point-distance matching
failure: the tray introduced in the experimental baseline
(Section~\ref{sec:baseline}), a target zone tens of centimeters wide,
accepts any placement inside its bounds.} A \rev{placement} that
exactly echoes a target's registered coordinates is treated as
intentional and left untouched. One that merely drifted into a zone's
footprint is nudged back out along the axis needing the smaller
correction. The nudge only fires when the zone's surface sits
meaningfully above the bare table; otherwise the near-table-height
generic ``table'' zone, which legitimately contains most of the
workspace, would push objects off the table's edge.

\subsection{Width-aware tool offset}
\label{sec:fix-tcp}

Because the gripper's fingers pivot, the flange-to-contact distance
depends on the commanded opening width. \harness{} measures this
distance at three representative widths: fully closed, mid-open, and
fully open. At plan time, it uses piecewise-linear interpolation
between these measurements and sets the executor's tool transform for
each pick according to the planned grasp width. Pick descent is
additionally capped so that tall objects are gripped nearer their top,
keeping the gripper's rigid mounting body clear. This \rev{resolves}
the width-dependent tool offset failure.

\subsection{Best-effort grasp-width estimation}
\label{sec:fix-width}

In the VLM \rev{implementation}, grasp width is estimated by
deprojecting the left and right edges of a detected bounding box.
Sampling exactly on the silhouette boundary is the worst location for
valid depth, and one failed sample used to abort the whole plan, over
a parameter that does not even affect whether pick and release are
safe. \harness{} samples slightly inset from the silhouette edge, and
isolates width estimation into its own best-effort call that degrades
to the model's estimate on failure instead of aborting. This
\rev{resolves} the silhouette-edge sampling failure, without yet
making width depth-verified in every case, a limitation
Section~\ref{sec:discussion} returns to.

\subsection{Zone-aware spacing: a partial mitigation}
\label{sec:zone-spacing}

For a zone that receives multiple \rev{placements} in one plan,
\harness{} lays the occupants out on a grid centered on the zone's
anchor and clamped to its registered bounds, falling back to an
unbounded grid only when no known zone matches. This \rev{addresses}
the \emph{static} half of the sequential-placement collision: final
positions no longer overlap, nor overflow a zone smaller than the
model's unclamped guess. The \emph{dynamic} half remains
\rev{unresolved}. The approach trajectory is not checked against
already-placed objects, and a released object's residual motion is not
constrained at all. The evaluation confirms this as the dominant
residual failure.

\subsection{Validation of the corrections}
\label{sec:validation}

Each correction was developed against real, previously failing
production runs, not synthetic cases. The baseline's per-run debug
artifacts let us replay the exact inputs that produced a given
hardware failure and confirm that the corrected pipeline now produces
the hand-derived expected geometry for the same input. This
retrospective replay, together with a unit-level regression suite
(45/45 passing across all reasoning-strategy implementations at the
time of writing), gives us confidence that each correction addresses
its diagnosed root cause. What remains is to measure the harness's
effect end to end, which is the subject of the next section.
